\scanpage{216}\mspage{161}

% MS 1–2
sur $S$), sous réserve que $W$ satisfait l'axiome L \unsure{5|6} (cf page 43).

% MS 3
2) Soit

% MS 4
\[
 \unsure{\xi|\zeta}'\deleted{\unsure{\in}}=\mathrm{hot}_{S',\unsure{W|w}}^{\mathrm{lc}}(X')\quad\text{dans }\mathrm{HOT}_W(S'),
\]

% MS 5
on a $f_!^0(\xi')=\mathrm{hot}_{S,W}^{\mathrm{lc}}(X')$,

% MS 6–9
mais pour expliciter $f_0^*$ \deleted{\unsure{$\xi$}} de cet élément, comme $X'$ n'a aucune raison d'être \deleted{\unsure{fibré} fibré} un \unsure{HA|H$\Lambda$}-fibré sur $S$ (auquel cas il suffirait de prendre $f^*(X')\overset{\mathrm{def}}=X'\times_S S'$), il faut d'abord prendre $C_S(X')$, puis appliquer $f^*$ :

% MS 10
\[
 f_0^*f_!^0(X')=\deleted{\unsure{C_S(X')}}\mathrm{hot}_{S',W}^{\mathrm{lc}}(C_S(X')\times_S S')
\]

% MS 11
Il faut donc définir

% MS 12
\[
 X'\xrightarrow{\beta_{0\,X'}} C_S(X')\times_S S'
 \qquad\text{dans }(\mathrm{Cat}_{/S'})W_{S'}^{-1}.
\]

% MS 13–15
Or on trouve une flèche canonique dans $\mathrm{Cat}_{/S'}$, déduite du $S$-morphisme $i_{X'}:X'\unsure{\text{---}|\longrightarrow} C_S(X')$ i.e. du carré comm. en traits pleins

% MS 16–18
\diagram{\begin{tikzcd}[ampersand replacement=\&,column sep=small,row sep=small]
 X'\arrow[rr,no head,"i_{X'}"]\arrow[dd,"p_{X'}"']\arrow[dr,dashed,"{\beta_{0\,X'}}"'] \& \& C_S(X')\arrow[dd]\\
 \& C_S(X')\times_S S'\arrow[ur,dashed,"{\mathrm{pr}_1}"']\arrow[dl,dashed]\arrow[d,phantom,"{\text{\unsure{comm.}}}"] \&\\
 S'\arrow[rr,"f"] \& \& S
\end{tikzcd}}

% MS punctuation at right of square
,

% MS 19–21
comme l'unique flèche pointillée $\beta_{0\,X'}$ qui soit un $S'$-morphisme (triangle gauche comm.) et qui rende commutatif le triangle supérieur.

% MS 22–25
Si $f\in W$, alors $\mathrm{pr}_1$, déduit de $f$ par chgt. de base $C_S(X')\longrightarrow S$ qui est parfait, est dans $W$ par L 5, on sait qu'il en est de $=$ de $i_{X'}$, donc aussi de $\beta_{0\,X'}$, on gagne !

% archive sheet number: 215
